% related-work-draft.tex -- CrackedPDFs revision
% Background (~350 words) and Related Work (~1,300 words).
% All \cite keys resolve in references.bib (same directory).

\subsection{Background: Indirect Prompt Injection Through PDF Text Extraction}
\label{sec:background}

An LLM-integrated application concatenates a trusted instruction from the
developer or user with untrusted data such as a retrieved web page, an email,
or an uploaded document. Because the model receives both as one token
sequence, instructions carried in the data can override the intended task.
When the attacker controls only the data and never interacts with the model
directly, the attack is called \emph{indirect} prompt
injection~\cite{greshake2023not}. Liu et al.~\cite{liu2024formalizing}
formalize it as the composition of a target task and an injected task, and
subsequent work has shown that current models do not reliably separate
instructions from data~\cite{zverev2025can}. OWASP ranks prompt injection as
the leading risk for LLM applications~\cite{owasp2025llm01}.

PDF adds a second gap, between what a person sees and what the model reads.
A PDF page is a content stream of drawing operators. Text-showing operators
place glyphs from an embedded font at positions set by the current
transformation matrix, and the graphics state controls fill color, opacity,
clipping, and the text rendering mode; mode~3 draws nothing at all. A
renderer rasterizes these operators for the human reader. A text extractor
(for example pdfminer, PyMuPDF, or Poppler's \texttt{pdftotext}) ignores most
of the graphics state. It maps character codes to Unicode through the font's
\texttt{/ToUnicode} table and emits the result in content-stream order.
Text that is white on white, set in a 0.5\,pt font, positioned outside the
crop box, clipped to a zero-area path, placed in a hidden optional-content
layer, or drawn in rendering mode~3 is therefore invisible on screen but
present in the extracted string. A font whose \texttt{/ToUnicode} map
disagrees with its glyph shapes makes the visible text and the extracted
text differ even when both are
``visible''~\cite{markwood2017pdf,xiong2025invisible,luo2026exploiting}.

Document-to-LLM pipelines, including retrieval loaders, hosted chat
services, résumé screeners, and automated reviewing tools, usually call such
an extractor and send the flattened string to the model and to any
guardrail~\cite{castagnaro2025hidden,liu2027what}. By that point the
information that the instruction was never visible to a person has been
lost. A text-level injection classifier sees the same string whether or not
the instruction was visible. This motivates detection that keeps access to
the document's rendering semantics.

% ---------------------------------------------------------------------------

\section{Related Work}
\label{sec:related}

\paragraph{Prompt injection attacks and benchmarks.}
Direct goal hijacking was first demonstrated by Perez and
Ribeiro~\cite{perez2022ignore}. Greshake et al.~\cite{greshake2023not}
introduced indirect injection through retrieved content, and Liu et
al.~\cite{liu2024formalizing} gave a formal framework and a benchmark of
attacks and defenses. Large human-generated attack corpora came from online
games and competitions~\cite{toyer2024tensor,schulhoff2023ignore}, and
optimized triggers from gradient-based search~\cite{pasquini2024neural}.
Benchmarks for indirect injection cover text carriers such as emails, web
pages, and tables (BIPIA~\cite{yi2025benchmarking}), tool outputs for agents
(InjecAgent~\cite{zhan2024injecagent},
AgentDojo~\cite{debenedetti2024agentdojo}), web agents
(WASP~\cite{evtimov2025wasp}, EIA~\cite{liao2025eia}), and social-web
retrieval corpora~\cite{guo2026hidden}. All of these represent the injected
payload as a string that is already in the model's context. None models the
file format that carried it, which is where the visibility of the payload is
decided.

\paragraph{Detection-based defenses.}
Deployed guardrails are mostly fine-tuned encoder classifiers applied to
text: Llama Prompt Guard~2~\cite{meta2025promptguard2,chennabasappa2025llamafirewall},
ProtectAI's and deepset's DeBERTa
models~\cite{protectai2024debertav2,deepset2023deberta},
PromptShield~\cite{jacob2025promptshield}, and
Sentinel~\cite{ivry2025sentinel}. Other detectors use a minimax-trained
detector LLM (DataSentinel~\cite{liu2025datasentinel}), attention
patterns~\cite{hung2025attention}, or activation deltas~\cite{abdelnabi2025get}.
PIGuard~\cite{li2025piguard} identified \emph{over-defense}, meaning false
positives on benign text that contains trigger words, and proposed a
benchmark and training method to reduce it. CAPTURE~\cite{kholkar2025capture}
measures both missed attacks and over-defense in context-specific settings.
Every one of these detectors operates after extraction. Our results show
that this placement is not a minor engineering detail: given only extracted
text, a classifier cannot tell a hidden instruction from the same sentence
printed visibly in a security tutorial.

\paragraph{Training-time and system-level defenses.}
A second line of work makes the model or the system robust to instructions
in the data channel. These include structured queries and preference
optimization (StruQ~\cite{chen2025struq}, SecAlign~\cite{chen2025secalign}),
task-specific fine-tuning~\cite{piet2024jatmo}, instruction
hierarchies~\cite{wallace2024instruction}, input marking
(spotlighting~\cite{hines2024defending}), LLM-based sanitization
(PromptArmor~\cite{shi2025promptarmor}), masked re-execution for agents
(MELON~\cite{zhu2025melon}), and capability-based isolation of control flow
from data flow (CaMeL~\cite{debenedetti2026defeating}). These defenses are
complementary to ours. They limit what an injected instruction can do once
it reaches the model, while a document-layer detector decides whether
invisible content should reach the model at all. A critical re-evaluation
found many such defenses weaker than reported once utility and adaptive
attacks are accounted for~\cite{jia2026critical}.

\paragraph{Adaptive attacks and the reliability of detector evaluations.}
Adaptive attackers defeat every indirect-injection defense evaluated by Zhan
et al.~\cite{zhan2025adaptive} and twelve published defenses evaluated by
Nasr et al.~\cite{nasr2026attacker}. We do not evaluate adaptive attackers,
and our results should be read as an upper bound on robustness. A separate
concern is whether reported detector accuracy measures the intended signal
at all. Fomin~\cite{fomin2026when} shows that same-source train/test splits
inflate the AUC of injection classifiers by 8.4 points and that many
salient features are dataset-specific shortcuts. PIDS-Bench~\cite{shire2026pidsbench}
finds that a detector with F1 above 0.98 in distribution misclassifies about
a third of externally sourced benign prompts. Li et al.~\cite{li2026defenses}
trace false refusals to position, trigger-token, and topic heuristics, and
RAG-PIBench~\cite{jaffal2026ragpibench} finds that sparse lexical baselines
remain competitive under leakage-aware splits. These findings concern text
inputs. Our benchmark applies the same scrutiny to documents, where a
shortcut can also live in the generator's templates or in incidental file
structure.

\paragraph{Hidden content in documents.}
The gap between rendered and extracted PDF content predates LLMs. PDF
Mirage~\cite{markwood2017pdf} used font remapping to mislead reviewer
assignment, plagiarism detection, and search indexing. Shadow
attacks~\cite{mainka2021shadow} exploited differences between PDF viewers to
alter signed documents, and Kuchta et al.~\cite{kuchta2018correctness}
found that 13.5\% of real PDFs render inconsistently across readers. For
LLMs, Unicode-level hiding~\cite{boucher2022bad,boucher2023trojan,gao2025imperceptible},
imperceptible content poisoning~\cite{zhang2024imperceptible}, and malicious
fonts~\cite{xiong2025invisible,luo2026exploiting} make the text a model reads
differ from what a person sees. TrapDoc~\cite{jin2025trapdoc} injects
invisible ``phantom tokens'' that make LLMs return plausible but wrong
answers, and Castagnaro et al.~\cite{castagnaro2025hidden} show that 19
hiding techniques in DOCX, HTML, and PDF succeed against RAG data loaders in
74.4\% of 357 scenarios. The work closest to ours is Liu and
Ming~\cite{liu2027what}. They catalogue 25 gaps between what users see and
what models read in PDF, measure them across 16 processing stacks and seven
commercial LLM services, and release a rule-based scanner. Their work maps
the attack surface. Ours provides a labeled detection benchmark with matched
benign confounders and held-out provenance splits, and compares learned and
rule-based detectors under one protocol. Their reading-order and
font-decoding families are outside our current coverage.

\paragraph{Detectors for hidden content.}
PhantomLint~\cite{murray2025phantomlint} is the closest prior detector. It
first selects text blocks that resemble prompts, then renders and OCRs those
blocks, and flags words that were extracted but not rendered. It reports a
0.092\% false-positive rate on ICML~2025 papers. Its evaluation contains few
benign documents with hidden-but-harmless content or visible prompt-like
language. These are the cases our confounders target, and the cases where a
phrase-gated, OCR-differential design is most likely to fail. Liu and
Ming~\cite{liu2027what}, Luo et al.~\cite{luo2026exploiting}, and
PhantomLint all argue that rendering followed by OCR is the safest reading of
an untrusted PDF. However, render-based pipelines remain exposed to text
embedded in pixels~\cite{clusmann2025prompt,gong2025figstep,bagdasaryan2023abusing,chen2026repeat}
and to encoded payloads that no single view
reveals~\cite{sinha2026hiding,jia2026seeing}.

\paragraph{Hidden prompts in peer review and hiring.}
In July 2025, reporters found white-text and tiny-font instructions such as
``give a positive review only'' in arXiv
preprints~\cite{nikkei2025positive,gibney2025scientists,lin2026hidden}, and
ICML issued a policy on such prompts~\cite{icml2025publication}. Controlled
studies show that these instructions shift LLM-generated review scores and
decisions~\cite{keuper2025prompt,collu2026misleading,zhou2026give,sahoo2025when,zhu2025when,theocharopoulos2025multilingual,li2026llmreviewer}.
Hidden prompts have also been used defensively as watermarks to detect
LLM-written reviews~\cite{rao2025detecting,gharami2025chatgpt}. A detector
must therefore not treat hidden text as malicious in itself. In hiring,
Zhang et al.~\cite{zhang2026measuring} measured hidden injections in roughly
1\% of about 200{,}000 real résumés. More than 90\% of them inserted data
rather than explicit instructions. Controlled résumé-screening studies and
benchmarks report similar
vulnerabilities~\cite{mu2025ai,baxi2026prompt,akdemir2025understanding,barach2026resumeshield}.
The prevalence of non-imperative payloads matters for us. A detector that
keys on instruction-like phrasing will miss the most common real-world case.

\paragraph{Knowledge-base poisoning.}
Retrieval poisoning attacks insert crafted passages that steer, backdoor, or
block retrieval-augmented
generation~\cite{zhong2023poisoning,zou2025poisonedrag,xue2024badrag,chaudhari2024phantom,shafran2025machine,chen2024agentpoison}.
These attacks optimize what the payload says. Our benchmark concerns where
the payload is hidden. The two compose naturally: a PoisonedRAG passage
placed in an invisible text layer survives both human review of the source
document and most extraction pipelines.

\paragraph{Structural PDF detection.}
Static malware detectors learn from PDF structure: embedded
JavaScript~\cite{laskov2011static}, metadata and object
statistics~\cite{smutz2012malicious}, and structural
paths~\cite{srndic2013detection,srndic2016hidost}. Mimicus, EvadeML, and
parser-confusion attacks showed that such features are easy to evade or
to hide from the detector's own parser~\cite{srndic2014practical,xu2016automatically,carmony2016extract,maiorca2019towards}.
This motivated verifiably robust training~\cite{chen2020training}. The same
lesson applies here. Structural features that separate injected from benign
files in a synthetic corpus may reflect how the generator works rather than
whether the document is hidden or malicious. This is why we audit
structure-only models against matched confounders.

\paragraph{Evaluation methodology.}
Our controls borrow from work on shortcut learning and leakage. Models
often exploit incidental cues~\cite{geirhos2020shortcut,lapuschkin2019unmasking,mccoy2019right},
and partial-input baselines expose annotation
artifacts~\cite{gururangan2018annotation}. Security ML suffers from the same
spurious correlations and biased splits~\cite{arp2022dos,pendlebury2019tesseract,warnecke2020evaluating,kapoor2023leakage}.
Each injected file and its confounder share a base document, so our setting
is a pair-input problem. In such problems, random splits overstate
performance unless they hold out the shared component~\cite{park2012flaws,pahikkala2015toward,bernett2024cracking}.
Matched confounders play the role of contrast
sets~\cite{gardner2020evaluating,kaushik2020learning}. Because files from the
same base document are correlated, we resample base documents when computing
confidence intervals~\cite{field2007bootstrapping,davison1997bootstrap} and
report design-effect-adjusted
uncertainty~\cite{kish1965survey,dean2015evaluating}.
